\documentclass[10pt]{article}

\usepackage{amsmath}
\usepackage{amssymb}
\usepackage{amsthm}
\usepackage{ mathrsfs }
\usepackage{ mathtools}
\usepackage{enumerate}
\usepackage{bbm}
\usepackage{lipsum}
\usepackage{fancyhdr}
\usepackage{tikz-cd} 
\usepackage{adjustbox} 
\usetikzlibrary{arrows}
\newcommand{\midarrow}{\tikz \draw[-triangle 90] (0,0) -- +(.1,0);}
\tikzset{commutative diagrams/.cd,
mysymbol/.style={start anchor=center,end anchor=center,draw=none}
}
\newcommand\comsymb[2][\circlearrowleft]{%
  \arrow[mysymbol]{#2}[description]{#1}}

\newtheorem{theorem}{Theorem} [section] 
\newtheorem{proposition}{Proposition}[section] 
\newtheorem{definition}{Definition}[section] 
\newtheorem{lemma}{Lemma}[section] 
\newtheorem{notation}{Notation}[section] 
\newtheorem{remark}{Remark}[section] 
\newtheorem{corollary}{Corollary} [section] 
\newtheorem{terminology}{Terminology}[section] 
\newtheorem{fact}{Fact}[section] 
\newtheorem{example}{Example}[section] 
\newtheorem{claim}{Claim}
\newenvironment{claimproof}[1]{\par\noindent\underline{Proof:}\space#1}{\hfill $\blacksquare$}

\DeclareMathOperator{\tmod}{mod}
\DeclareMathOperator{\triv}{triv}
\DeclareMathOperator{\ind}{ind}
\DeclareMathOperator{\straw}{s.t.}
\DeclareMathOperator{\ev}{ev}
\DeclareMathOperator{\pr}{pr}
\DeclareMathOperator{\Aut}{Aut}
\DeclareMathOperator{\Map}{Map}
\DeclareMathOperator{\Stab}{Stab}
\DeclareMathOperator{\Ind}{Ind}
\DeclareMathOperator{\GL}{GL}
\DeclareMathOperator{\SL}{SL}
\DeclareMathOperator{\SO}{SO}
\DeclareMathOperator{\Sp}{Sp}
\DeclareMathOperator{\esssup}{esssup}
\DeclareMathOperator{\abs}{abs}
\DeclareMathOperator{\rel}{rel}
\DeclareMathOperator{\diam}{diam}
\DeclareMathOperator{\rank}{rank}
\DeclareMathOperator{\Hom}{Hom}
\DeclareMathOperator{\Homk}{Hom_{k-alg}}
\DeclareMathOperator{\Ker}{Ker}
\DeclareMathOperator{\Image}{Im}
\DeclareMathOperator{\Dom}{Dom}
\DeclareMathOperator{\grad}{grad}
\DeclareMathOperator{\divergence}{div}
\DeclareMathOperator{\rk}{rk}
\DeclareMathOperator{\Span}{Span}
\DeclareMathOperator{\mSpec}{mSpec}
\DeclareMathOperator{\MaxSpec}{MaxSpec}
\DeclareMathOperator{\interior}{int}
\DeclareMathOperator{\supp}{supp}
\DeclareMathOperator{\id}{id}
\DeclareMathOperator{\sgn}{sgn}
\DeclareMathOperator{\Mat}{Mat}
\DeclareMathOperator{\PD}{PD}
\DeclareMathOperator{\ext}{ext}
\DeclareMathOperator{\vol}{vol}
\DeclareMathOperator{\inte}{int}
\DeclareMathOperator{\diverg}{div}
\DeclareMathOperator{\curl}{curl}
\DeclareMathOperator{\image}{im}
\DeclareMathOperator{\dR}{dR}
\DeclareMathOperator{\Ob}{Ob}
\DeclareMathOperator{\Mnfd}{Mnfd}
\DeclareMathOperator{\OpEmb}{OpEmb}
\DeclareMathOperator{\Spec}{Spec}
\DeclareMathOperator{\Spa}{Spa}
\DeclareMathOperator{\Sper}{Sper}
\DeclareMathOperator{\op}{op}
\DeclareMathOperator{\con}{con}
\DeclareMathOperator{\Gal}{Gal}
%\newcommand*{\name}[\num_arguments][default values]{{\color{#1}\Large #2}}
\newcommand{\defeq}{\vcentcolon=}
\newcommand{\Ouv}{\mathcal{O}}
\newcommand{\norm}[1]{\Vert #1 \Vert}
\newcommand{\opNorm}[2]{\norm{#1}_{#2\to#2}}
\newcommand{\normL}[3]{\norm{#1}_{L^{#2}(#3)}}
\newcommand{\N}[0]{\mathbb{N}}
\newcommand{\m}[0]{\mathfrak{m}}
\newcommand{\R}[0]{\mathbb{R}}
\newcommand{\Z}[0]{\mathbb{Z}}
\newcommand{\C}[0]{\mathbb{C}}
\newcommand{\Q}[0]{\mathbb{Q}}
\newcommand{\Qc}[0]{\mathfrak{Q}_C}
\newcommand{\F}[0]{\mathbb{F}}
\newcommand{\G}[0]{\mathbb{G}}
\newcommand{\A}[0]{\mathbb{A}}
\newcommand{\T}[0]{\mathbb{T}}
\newcommand{\Para}[0]{\mathbb{P}}
\newcommand{\M}[0]{\mathcal{M}}
\newcommand{\maniN}[0]{\mathcal{N}}
\newcommand{\cinf}[0]{\mathcal{C}^\infty}
\newcommand{\torus}[0]{\mathbb{T}}
\newcommand{\sep}[0]{\:|\:}
\newcommand{\pd}[2]{{\frac {\partial #1} {\partial #2}}}
\newcommand{\compinc}[0]{\subset \subset}
\newcommand{\sH}[0]{\mathscr{H}}
\newcommand{\ab}[1]{\vert #1 \vert}
\newcommand{\st}[0]{\:\straw\:}
\newcommand{\g}[0]{\mathfrak{g}}
\newcommand{\p}[0]{\mathfrak{p}}
\newcommand{\U}[0]{\mathfrak{U}}
\newcommand{\fB}[0]{\mathfrak{B}}
\newcommand{\fib}[1]{%
  \mathbin{\mathop{\times}\limits_{#1}}%
}
\newcommand{\tens}[1]{%
  \mathbin{\mathop{\otimes}\displaylimits_{#1}}%
}

\title{Rigid Analytic Geometry Week 8 Solution}
\author{So Murata}
\date{SoSe 26, University of Bonn}


\begin{document}
\maketitle

\par\textbf{Problem 1}
By Problem 8 of the previous sheet, there is $R\in (r,1]$ such that for $r<\ab{x}\leq R$, we have,
\begin{equation*}
    \ab{f_i(x)}=n_i\left({\frac {\ab{x}} {r}}\right)^{j_i}.
\end{equation*}
Since the ordering is lexicographic, we have two cases for each $i=1,\cdots,m$, such that
\begin{enumerate}
    \item $n_i<n_0$
    \item $n_i=n_0$ and $j_i\leq j_0$.
\end{enumerate}
For the first case, 
\begin{equation*}
    n_i\left({\frac R r}\right)^{j_i}<n_0\left({\frac R r}\right)^{j_0}.
\end{equation*}
holds by taking suitably small $R$ and using the continuity of $\ab{\cdot}$, we obtain the inequality.
\par For the second case, $j_i\leq j_0$ implies 
\begin{equation*}
    \left({\frac {\ab{x}} {r}}\right)^{j_i}<\left({\frac {\ab{x}} {r}}\right)^{j_0}.
\end{equation*}
This shows that 
\begin{equation*}
    r<\ab{x}\leq R\Rightarrow x\in\mathcal{R}(f_0|f_1,\cdots,f_m).
\end{equation*}
This proves the first part.
\par For the second part, we again have two cases,
\begin{enumerate}
    \item $n_i>n_0$
    \item $n_i=n_0,j_i>j_0$.
\end{enumerate}
For the first case, again by the continuity of $\ab{\cdot}$, we have,
\begin{equation*}
    n_i\left({\frac R r}\right)^{j_i}>n_0\left({\frac R r}\right)^{j_0}.
\end{equation*}
And similarly for the other case. Therefore,
\begin{equation*}
    \vert f_i(x)\vert>\ab{f_0(x)}
\end{equation*}
for all $r<\ab{x}\leq R$. Therefore,
\begin{equation*}
    \Omega\cap B_{\leq R}(0)\subseteq B_{\leq r}(0).
\end{equation*}
\par\textbf{Problem 5}\\
By Problem 4, we can take $F=\{x_1,\cdots,x_r\}\subset \mathcal{S}(B)$ such that 
\begin{equation*}
    z\in\bigcap_{j=1}^r S_R(x_j)\ni Z\Rightarrow \ab{f_0(z)}=\norm{f_0}_B.
\end{equation*}
Then by the definition of norms and the assumption,
\begin{equation*}
    \ab{f_i(x)}\leq \norm{f_i}_B\leq \norm{f_0}_B=\ab{f_0(0)}, 1\leq i\leq m.
\end{equation*}
This shows that 
\begin{equation*}
    \bigcap_{S\in F}S\subseteq \Omega.
\end{equation*}
\par The second part is similarly shown by applying Problem 4 to $f_i$ such that 
\begin{equation*}
    \norm{f_0}_B<\norm{f_i}_B.
\end{equation*}
That is 
\begin{equation*}
    x\in\bigcap_{S\in F}S\Rightarrow \ab{f_0(x)}<\ab{f_i(x)}.
\end{equation*}
Therefore,
\begin{equation*}
    \Omega\cap\bigcap_{S\in F}S=\emptyset.
\end{equation*}
\par\textbf{Problem 7}\\
Set,
\begin{equation*}
    \Phi:A\times M\to B\times M, (a,m)\mapsto(a\cdot1_B,m).
\end{equation*}
Since $a\mapsto f(a)$ is continuous $\Phi$ is continuous. By assumption $\rho:B\times M\to M, (b,m)\mapsto bm$ is continuous. Therefore 
\begin{equation*}
    \rho\circ\Phi:A\times M\to M
\end{equation*}
is continuous. This shows $M$ is an adic $A$-module.
\par\textbf{Tutor's solutions}
\par\textbf{Problem 1,2}
Accorind to Sheet 9 Problem 8, one can replace there is $r_+,r_-,0<r_-<r<r_+$ in Sheet 7 Problem 1,3 by there is $R, r<R<1$ and the proof is the same.
\par\textbf{Problem 3}
We have,
\begin{equation*}
    B = \bigcap_{S\in\mathcal{S}}S \not=\emptyset,
\end{equation*}
where $B\in\xi$. if $S,S'$ are different elements of $B=K^\circ$, Take $S\not=S'\in\mathcal{S}(B)$ Say $S=S_1(x),S' = S_1(y)$. Then $\ab{x-y}=1$, else $S=S'$.
For any $z\in B,\ab{z-x}<1$ and $\ab{z-y}<1$, cannot happen at the same time. Therefore $z\in S$ or $z\in S'$. Thus $S\cup S'=B$. $B\in\xi\Rightarrow$ one of $S,S'$ is in $\xi$. 
\begin{remark}
    Any finite intersection of elements in $\mathcal{S}(B)$ is non-empty.
\end{remark}
\begin{proof}
    Assume $B=K^\circ$. Take$ x\in K^\circ$ Then $\ab{y-x}=1$ if and only if $y-x\not\in K^{\circ\circ}$. Take $x_1,\cdots,x_m\in K^\circ$. Then,
    \begin{equation*}
        \bigcap_{i}\mathcal{S}_1(x_i) = \{y\in K^\circ\sep \overline{y} \not=\overline{x}_1,\cdots,\overline{x}_m\text{ in }k = K^\circ/K^{\circ\circ}\}.
    \end{equation*}
    Note that $K$ is algebraically closed, so $\ab{k} = \infty$. Therefore $\bigcap_i S_1(x_i) \not=\emptyset$.
\end{proof}
\par\textbf{Problem 4}
Assume $\norm{g|\T_1} = 1$. Then $\ab{g(x)}<\norm{g|\T_1}$ if and only if $g(x)\in K^{\circ\circ}$ if andon ly if $\overline{x}$ is a root of $\overline{g}$ in $k = K^\circ/K^{\circ\circ}$. Note that $\overline{g}$ is a polynomial in $k$, so we may take all of its roots $\alpha_1,\cdots,\alpha_n$. Choose $x_i$ to be a lift of $\alpha_i$, then 
\begin{equation*}
    \ab{g(x)}<\norm{g|\T_1}\iff x-x_i\in K^{\circ\circ}\iff x\in\bigcup_i(x_i+K^{\circ\circ}).
\end{equation*}
Take the complement.
\par\textbf{Problem 5}
Take a finite set $RF\subseteq\mathcal{S}(B)$ such that for $x\in\bigcap_{S\in F}S$,
\begin{equation*}
    \ab{f_i(x)}=\norm{f_i|\T_1}.
\end{equation*}
Then the conclusion is immediate.
\par\textbf{Problem 6}
For a rational open $\Omega$, $\Omega\in\xi_B\Rightarrow \exists F\subseteq \mathcal{S}(B)$ finite with 
\begin{equation*}
    \bigcap_{S\in F}F\subseteq\Omega.
\end{equation*}
Since each $S$ is in $\xi$, so does their finite intersection. Hence $\Omega\in\xi$. 
\par $\Omega\not\in\xi_B$, then there is a finite set $F\subseteq\mathcal{S}(B)$ such tat 
\begin{equation*}
    \Omega\cap\bigcap_{S\in F}S = \emptyset. 
\end{equation*}
Note that $\bigcap_{S\in F}S\in\xi$ therefore $\Omega\not\in\xi$. Thus $\xi=\xi_B$.
\par\textbf{Problem 7}
\begin{lemma}
    If $A$ is a Huber ring, $(A^\#,I)$ be a pair of definition, $C$ be an open subring of $A$, then 
    \begin{equation*}
        (C\cap A^\#,I^NA^\#)
    \end{equation*}
    is also a pair of definition.
\end{lemma}
\begin{proof}
    Since $C\cap A^\#$ is open, there is $N$ such that $I^N A^\#\subseteq C\cap A^\#$. Then $I^NA^\#$ is finitely generated $C\cap A^\#$-ideal. First 
    \begin{equation*}
        (I^NA^\#)^n(C\cap A^\#)
    \end{equation*}
    is open.
\end{proof}
Moreover, $I^NA^\#$ is finitely generated, also $(I^NA^\#)^n(C\cap A^\#)$ gives $0$-neighbourhood of $A$.
\par\textbf{Problem 7}
\begin{lemma}
    If $M$ is an adic $A$-module with respect to $(A^\#,I)$. Let $(A',J)$ be another pair of definition. Then $M$ is adic with respect to $(A',J)$.
\end{lemma}
\begin{proof}
    By the assumtpion, there is an open submodule $M^\#\subseteq M$ such that $I^n M^\#$ is a basis of $0$-neighbourhood. Let $M'=A'M^\#$. Then $M'$ is an open sub-$A'$-module of $M$. Now since $(A^\#,I)$ and $(A',J)$ define the same topology. $\{I^n A^\#\},\{J^n A'\}$ are cofinal in the sense of covering $I^nM^\#$, there is $N$ such that 
    \begin{equation*}
        J^NA'\subseteq I^nA^\#.
    \end{equation*} 
    Therefore $J^nM' = J^NA'M^\#\subseteq I^nA^\#M^\#=I^nM^\#$. Given $I^nM'$ there is $N$ such that $I^nA^\#\subsetrq J^nA'$.
    THerefore $I^NA^\#M^\#\subseteq J^nA'M^\#=J^nM'$. Thus $\{J^nM'\}$ is also a base for $0$-neighbourhood.
\end{proof}
\begin{theorem}
    \begin{equation*}
        (\Sp A)^\ast\leftrightarrow\Spa(A,A^\circ).
    \end{equation*}
    We now consider $A=\T_1$. 
    \begin{enumerate}
        \item A classical point $x_a$ consists of an element $a\in K^\circ$ and valuatio n$v_0(f) = \ab{f(a)}$.
        \item Classical disk points $x_{a,r}$ consits of an element $a\in K^\circ$ and $0<r\leq 1$ and the valuation,
        \begin{equation}
            f=\sum_{i}f_i(T-a)^i\mapsto \max_{i}\ab{f_i}r^i.
        \end{equation}
    \end{enumerate}
    This onlyl depends on $B_{\leq r}(a)$. 
\end{theorem}
\begin{proposition}
    For any $f\in\T_1$, $\ab{f(x)}$ is a constant for sufficently small $B_{r}\in\mathcal{M}$. Valuation 
\end{proposition}
\end{document}